\documentclass[11pt]{article}

\usepackage[margin=1.15in]{geometry}
\usepackage{amsmath,amssymb,amsthm}
\usepackage{booktabs}
\usepackage{microtype}
\usepackage{xcolor}
\usepackage[colorlinks=true,linkcolor=blue!50!black,citecolor=blue!50!black,urlcolor=blue!50!black]{hyperref}

\newtheorem{theorem}{Theorem}[section]
\newtheorem{lemma}[theorem]{Lemma}
\newtheorem{corollary}[theorem]{Corollary}
\newtheorem{proposition}[theorem]{Proposition}
\newtheorem{definition}[theorem]{Definition}
\newtheorem{remark}[theorem]{Remark}
\newtheorem{conjecture}[theorem]{Conjecture}

\newcommand{\N}{\mathbb{N}}
\newcommand{\Q}{\mathbb{Q}}
\newcommand{\Z}{\mathbb{Z}}
\newcommand{\Ztwo}{\Z_2}
\newcommand{\code}[1]{\texttt{#1}}
\newcommand{\jodd}[2]{j_{#1}(#2)}

% Lean-statement display: monospace with math symbols spliced in
\newenvironment{leanblock}
  {\begin{quote}\ttfamily\small\setlength{\parindent}{0pt}}
  {\end{quote}}

\title{The Collatz Conjecture at the Critical Line:\\
       Machine-Verified No-Go, Density, and Cycle Theorems}
\author{M.~Sharpe\thanks{\texttt{msharpe@irendi.com}. The mathematics,
  formalization, and text of this paper were produced in collaboration
  with Claude (Anthropic). All formal proofs are verified by the Lean~4
  kernel; see Section~\ref{sec:formal} for the audit.}}
\date{June 10, 2026}

\begin{document}
\maketitle

\begin{abstract}
We present a machine-verified account of where the Collatz conjecture
is provably hard and where it is provably almost true. On the negative
side, no \emph{finite-memory descent certificate} exists: there is no
potential $V(n)=n\,W(n)$ with $W$ bounded above and below away from
zero, no window $B$ and contraction factor $\rho<1$, such that every
large $n$ satisfies $V(T^t(n))\le\rho V(n)$ for some $t\le B$; the
obstruction is the $2$-adic fixed point $-1$, shadowed by the integers
$2^L-1$. On the positive side we formalize, in elementary integer
arithmetic throughout: two-sided \emph{master inequalities}
$3^j(n+1)-2^j \le 2^T\,T^T(n) \le 3^j(n+2^{T-j})-2^T$ for the Terras
map with $j$ odd steps among $T$, sharp on the shadowing orbit; the
\emph{density dichotomy} --- an orbit segment avoiding descent either
runs supercritical odd density $2^T<2\cdot 3^j$ or has exhausted the
$\log_2 n$-bit information budget of its start --- together with its
exact converse, so that non-descent profiles are pinned to an
$O(1/B)$-wide strip above the critical density $\log 2/\log 3$,
inhabited at every scale; the reduction of the full conjecture to
universal descent; Terras' parity-vector bijection and the binomial law
$\#\{r<2^k : j_k(r)=j\}=\binom{k}{j}$; Terras' theorem in finite form
--- the integers in $[2^k,2^{k+1})$ not descending within $k$ steps
number at most $3^k/2^{\lfloor 17(k-1)/27\rfloor}\approx 2^{0.955k}$
--- with no real analysis anywhere; and Diophantine cycle bounds
excluding, conditionally on the computational verification below
$2^{71}$, any nontrivial cycle with at most $183$ halvings, the finite
check performed by the Lean kernel itself. We further report a
systematic experimental program on certificates \emph{beyond} the no-go
theorem's hypotheses --- weighted automata reading all binary digits of
$n$ --- identifying the complete spectrum of positive-drift rational
cycles as adversaries and three mechanisms (injection, global rewrite,
shift re-alignment) that empirically refute every $1$- and $2$-state
digit certificate. All formal results are verified in Lean~4 over
\textsf{mathlib} with no \code{sorry} and no axioms beyond standard
foundations. The conjecture remains open; these results fence, with
sharp constants, exactly where it lives.
\end{abstract}

\section{Introduction}

The Collatz function $C(n) = n/2$ for even $n$ and $C(n) = 3n+1$ for
odd $n$ is conjectured to drive every positive integer to $1$. The
conjecture has been verified by exhaustive computation for all
$n \le 2^{68}$ in the published record \cite{barina}, with the
distributed effort since extending beyond $2^{71}$, and is the subject
of an extensive literature surveyed by Lagarias \cite{lagarias-survey,
lagarias-book}.

Because the conjecture is so easily verified on any finite range, a
perennially attractive proof strategy --- attempted in countless
preprints, and dissected in Section~\ref{sec:casestudy} --- runs as
follows:
\begin{enumerate}
\item partition the integers into finitely many classes (typically
  residues modulo $2^k$);
\item show by finite computation that every class guarantees a decrease
  in some bounded number of steps, perhaps as measured by a weighted
  potential;
\item conclude that every large $n$ descends below itself, and finish by
  strong induction into the computationally verified range.
\end{enumerate}
Step 3 is sound, and step 2 \emph{feels} like a finite computation. The
first contribution of this paper is a precise and machine-verified
proof that step 2 is impossible --- not merely unproven, but false for
every choice of class structure, potential, and window
(Theorem~\ref{thm:nogo}, \code{no\_finite\_certificate}).

The second contribution is the positive complement: a formalization of
exactly how close the impossible computation comes to being possible.
We prove two-sided, division-free \emph{master inequalities} for the
Terras map (Section~\ref{sec:corridor}), from which follow: the
\emph{density dichotomy} --- avoiding descent over a window of length
$T$ forces odd-step density above the critical value
$\log 2/\log 3 \approx 0.6309$, as long as the orbit is inside the
$\log_2 n$-bit information budget of its starting point --- and its
exact converse (Section~\ref{sec:critical}): integers exist, at every
scale and arbitrarily large, that avoid descent at densities within
$O(1/B)$ of critical. The non-descent profiles achievable over any
finite window thus form a strip of width $O(1/B)$ hugging the critical
line from above, \emph{inhabited at every scale}. Since the conjecture
is equivalent to universal descent
(Theorem~\ref{thm:reduction}, \code{collatz\_iff\_descent}), it is
exactly the assertion that no single integer tracks this vanishing
strip at every scale simultaneously --- an assertion the no-go theorem
shows no bounded-memory argument can decide.

The third contribution quantifies how rare the strip's inhabitants are.
We formalize Terras' parity-vector bijection --- the first $k$
parities of an orbit are determined by, and determine, $n \bmod 2^k$
--- with its corollaries that \emph{every} parity pattern is realized
and that exactly $\binom{k}{j}$ residues modulo $2^k$ take $j$ odd
steps (Section~\ref{sec:parity}). Combining the binomial law with the
dichotomy yields Terras' theorem \cite{terras} in finite form
(Section~\ref{sec:terras}): the integers in the dyadic window
$[2^k, 2^{k+1})$ that fail to descend within $k$ steps number at most
$3^k / 2^{\lfloor 17(k-1)/27 \rfloor} \approx 2^{0.955k}$, an
exponentially vanishing fraction $\approx 2^{-0.045k}$. The proof is
elementary throughout --- the tail estimate is the weighted binomial
theorem $2^m \sum_{j \ge m} \binom{k}{j} \le (1+2)^k$ in plain
integers, and the threshold $17/27$ comes from the single inequality
$3^{17} \le 2^{27}$ --- so the formalization involves no measure
theory and no real analysis. We believe this is the first
machine-verified density result for the Collatz map.

Fourth, cycles (Section~\ref{sec:cycles}). The two master inequalities
pinch on a cycle, giving the classical Diophantine constraint
$n(2^T - 3^j) < 3^j 2^{T-j}$ on every cycle element, fully formal; a
kernel-verified integer computation (plain \code{decide}; the lemma
depends on the axiom \code{propext} alone) then shows that,
conditional on the computational verification below $2^{71}$ ---
stated as an explicit hypothesis, since that computation is not
formalized --- no nontrivial cycle has $183$ or fewer halvings.

Finally, Section~\ref{sec:search} reports the experimental program
that books-ends the formal results. The linear-programming search over
joint $2$-adic/$3$-adic residue classes, whose systematic failure
motivated the no-go theorem, is complemented by two new instruments:
an enumeration of the \emph{positive-drift rational cycles} of the
Terras map --- the complete spectrum of shadowing adversaries, of
which the no-go theorem's $-1$ is only the extreme point --- and a
counterexample-guided (CEGIS) search over weighted automata that read
\emph{all} binary digits of $n$. Such automata evade the no-go
theorem's hypotheses, and one candidate found by the search is
genuinely instructive: a two-state automaton that tracks
$\nu_2(n+1)$ and thereby \emph{pre-pays} for every Mersenne rising
stretch. It and all its siblings nonetheless die to three identified
mechanisms --- injection, global rewrite, and shift re-alignment ---
that we conjecture kill every finite-state digit certificate.

The Collatz conjecture remains open, and nothing in this paper proves
otherwise. Section~\ref{sec:discussion} states exactly what the
results do \emph{not} rule out and where the residual difficulty
lives.

\section{Preliminaries}

For $n \in \N_{>0}$, the \emph{Collatz map} is
\[
C(n) = \begin{cases} n/2 & n \text{ even},\\ 3n+1 & n \text{ odd},\end{cases}
\qquad\text{and the \emph{Terras map} is}\qquad
T(n) = \begin{cases} n/2 & n \text{ even},\\ (3n+1)/2 & n \text{ odd}.\end{cases}
\]
The two generate the same orbits up to bookkeeping: one Terras step on
an odd $n$ equals two Collatz steps, since $3n+1$ is even. We work with
$T$ because every step performs exactly one halving, which makes size
estimates uniform; the formalization proves the translation
(\code{collatz\_two\_eq\_terras}) and states results for both maps.

For $T, n \in \N$ write
\[
\jodd{T}{n} \;=\; \#\bigl\{\, t < T : T^t(n) \text{ odd} \,\bigr\}
\]
for the number of odd steps among the first $T$ Terras steps
(\code{oddSteps}). The \emph{stopping time} of $n$ is the least $t$
with $T^t(n) < n$, if it exists. The conjecture is equivalent to:
every $n \ge 2$ has finite stopping time and the orbit of every $n$
reaches $1$; the first equivalence is made formal in
Theorem~\ref{thm:reduction}.

On the $2$-adic integers $\Ztwo$, the odd branch $n \mapsto (3n+1)/2$
extends continuously, and
\[
\frac{3 \cdot (-1) + 1}{2} = -1,
\]
so $-1$ is a fixed point of the odd branch, with multiplier $3/2 > 1$ in
archimedean size. No positive integer equals $-1$, but positive integers
approximate it $2$-adically to any precision: $2^L - 1 \equiv -1
\pmod{2^L}$. More generally every rational $a/q$ with $q$ odd lies in
$\Ztwo$ and is approximated by integers to any $2$-adic precision; this
becomes important in Section~\ref{sec:search}. The next section makes
the resulting shadowing exact for $-1$.

\section{The shadowing orbit}\label{sec:shadow}

\begin{lemma}[\code{terras\_mersenne}]\label{lem:shadow}
For all $L$ and all $j \le L$,
\[
T^j\!\left(2^L - 1\right) = 3^j\, 2^{L-j} - 1 .
\]
\end{lemma}

\begin{proof}
Induction on $j$. For $j = 0$ both sides equal $2^L - 1$. For the step,
let $m = 3^j 2^{L-j} - 1$ with $j < L$. Then $m = 2A - 1$ with
$A = 3^j 2^{L-j-1} \ge 1$, so $m$ is odd and
\[
T(m) = \frac{3(2A - 1) + 1}{2} = 3A - 1 = 3^{j+1} 2^{L-(j+1)} - 1.
\qedhere
\]
\end{proof}

In particular $T^j(2^L-1) \ge 2^L - 1$ for all $j \le L$
(\code{terras\_mersenne\_ge}), since $3^j 2^{L-j} \ge 2^L$; the orbit
rises strictly at every step until step $L$. Two immediate corollaries:

\begin{corollary}[\code{no\_uniform\_descent\_bound}]\label{cor:unbounded}
For every window $B$ and cutoff $N_0$ there exists $n \ge N_0$ with
$T^t(n) \ge n$ for all $t \le B$. (Take $n = 2^{B + N_0 + 1} - 1$.)
Equivalently, stopping times are unbounded.
\end{corollary}

\begin{corollary}[\code{uniform\_descent\_53\_false}]\label{cor:axiom}
The statement ``every $n > 300$ satisfies $C^t(n) < n$ for some
$t \le 53$'' is false: $n = 2^{28} - 1$ rises for its first $56$ Collatz
steps.
\end{corollary}

Corollary~\ref{cor:axiom} is stated in this specific form as a
concrete representative of the uniform-descent genre, calibrated so
that no counterexample exists below the threshold (Section~%
\ref{sec:casestudy} dissects how such claims come to be ``verified'').
Empirically the smallest violator is $n = 447 = 110111111_2$, whose six
trailing ones already exhibit the shadowing mechanism in miniature;
there are $1{,}933$ violators below $2 \times 10^5$.

\section{The no-go theorem}\label{sec:nogo}

\begin{definition}\label{def:cert}
A \emph{finite-memory descent certificate} for $T$ is a tuple
$(W, \lambda, \Lambda, B, N_0, \rho)$ with $W \colon \N \to \Q$,
$0 < \lambda \le W(n) \le \Lambda$ for all $n$, $B, N_0 \in \N$, and
$0 \le \rho < 1$, such that
\[
\forall\, n \ge N_0\;\; \exists\, t \in [1, B]:\quad
T^t(n)\; W\!\left(T^t(n)\right) \;\le\; \rho\; n\, W(n).
\]
\end{definition}

The definition is calibrated so that a certificate would prove the
conjecture:

\begin{proposition}\label{prop:motivation}
If a finite-memory descent certificate exists with $N_0$ at most the
computationally verified bound, the Collatz conjecture holds.
\end{proposition}

\begin{proof}[Proof sketch]
Set $V(n) = n\,W(n)$. Chaining certificate windows from any $n \ge N_0$
while the orbit remains $\ge N_0$ gives iterates $n_i$ with $V(n_i) \le
\rho^i V(n)$, hence $n_i \le (\Lambda/\lambda)\, \rho^i\, n \to 0$; so
the orbit must drop below $N_0$, where verified convergence to $1$
finishes the proof.
\end{proof}

The point of the multiplicative form is that it captures additive
log-potentials with bounded correction terms. If $S \colon \N \to Q$ is
\emph{any} assignment of states in a finite set $Q$ --- a residue
$n \bmod m$ for any modulus $m$ (in particular joint $2$-adic/$3$-adic
classes $m = 2^a 3^b$), the state of any finite automaton fed by the
orbit of $n$, any bounded window of past parities --- and
$\varphi \colon Q \to \mathbb{R}$ any weights, then
$V(n) = \log n + \varphi(S(n))$ with guaranteed decrease $\varepsilon >
0$ per window is the case $W = \exp \varphi \circ S$, $\rho =
e^{-\varepsilon}$, $\lambda = e^{\min \varphi}$, $\Lambda = e^{\max
\varphi}$ of Definition~\ref{def:cert}. The unweighted descent
arguments of Section~\ref{sec:casestudy} are the case $W \equiv 1$.

\begin{theorem}[\code{no\_finite\_certificate}]\label{thm:nogo}
No finite-memory descent certificate for $T$ exists.
\end{theorem}

\begin{proof}
Suppose $(W, \lambda, \Lambda, B, N_0, \rho)$ were one. From the bounds,
$0 < \lambda \le \Lambda$. Choose $k$ with $\Lambda \rho^k < \lambda$,
possible since $\rho < 1$. Let $L = kB + N_0 + 1$ and $n_0 = 2^L - 1
\ge N_0$.

We chain $k$ certificate windows inside the rising stretch of
Lemma~\ref{lem:shadow}. Inductively construct times $0 = \tau_0 <
\tau_1 < \dots < \tau_k$ with $\tau_i \le iB$ and
\begin{equation}\label{eq:chain}
T^{\tau_i}(n_0)\; W\!\left(T^{\tau_i}(n_0)\right) \;\le\; \rho^i\,
n_0\, W(n_0).
\end{equation}
For the step: $\tau_i \le iB \le kB < L$, so by
Lemma~\ref{lem:shadow} the current iterate $m_i = T^{\tau_i}(n_0)$
equals $3^{\tau_i} 2^{L - \tau_i} - 1 \ge n_0 \ge N_0$, and the
certificate applies to $m_i$, yielding $t \in [1,B]$ with
$T^{t}(m_i)\,W(T^t(m_i)) \le \rho\, m_i W(m_i)$; set $\tau_{i+1} =
\tau_i + t$ and combine with \eqref{eq:chain}.

At $i = k$: since $\tau_k \le kB < L$, the iterate $m_k =
T^{\tau_k}(n_0)$ still satisfies $m_k \ge n_0$, whence
\[
n_0\, \lambda \;\le\; m_k\, W(m_k) \;\le\; \rho^k\, n_0\, W(n_0)
\;\le\; n_0\, \Lambda\, \rho^k \;<\; n_0\, \lambda,
\]
a contradiction. (Note that the certificate hypothesis was used only on
integers of the single family $3^j 2^{L-j} - 1$; the theorem is thus
really a statement about the orbit of $-1 \in \Ztwo$.)
\end{proof}

\begin{corollary}[\code{no\_finite\_state\_certificate}]\label{cor:fsm}
For every finite type $Q$, every $S \colon \N \to Q$, and every
$w \colon Q \to \Q_{>0}$, there is no $(B, N_0, \rho)$ with $\rho < 1$
such that all $n \ge N_0$ admit $t \in [1,B]$ with $T^t(n)\,
w(S(T^t(n))) \le \rho\, n\, w(S(n))$.
\end{corollary}

\begin{proof}
Apply Theorem~\ref{thm:nogo} with $W = w \circ S$, $\lambda = \min_Q w$,
$\Lambda = \max_Q w$.
\end{proof}

\begin{remark}\label{rem:unbounded}
The boundedness of $W$ cannot be dropped: it is exactly what makes
contraction of $V$ assert something about descent of $n$
(Proposition~\ref{prop:motivation} uses both bounds). With unbounded
$W$ the quantity $V$ can contract while $n$ grows --- already
$W(n) = n^{-1-\varepsilon}$ gives a $V$ that contracts along every
rising stretch --- so ``certificates'' in that generality must carry a
validity floor $V \ge n^{\delta}$ to certify anything. Certificates
with unbounded $W$ satisfying such a floor are \emph{not} excluded by
Theorem~\ref{thm:nogo}; Section~\ref{sec:search} investigates them
experimentally. Likewise, if the window $B$ is allowed to depend on
$n$, the unweighted statement becomes the finiteness of stopping
times, i.e.\ the open descent half of the conjecture itself.
\end{remark}

\section{The master inequalities and the density corridor}
\label{sec:corridor}

Everything in this and the next four sections is elementary, exact, and
division-free; the formal proofs run in $\N$ with no truncated
subtraction. The engine is a pair of one-step identities: for odd $x$,
$2\,(T(x)+1) = 3(x+1)$ exactly, and for even $x$, $2\,(T(x)+1) = x+2$.
Induction along the orbit gives:

\begin{theorem}[master inequalities;
\code{terras\_growth\_bound}, \code{terras\_lower\_bound}]
\label{thm:master}
For all $T, n$, with $j = \jodd{T}{n}$,
\[
3^j (n+1) - 2^j \;\le\; 2^T\, T^T(n) \;\le\; 3^j\,(n + 2^{T-j}) - 2^T .
\]
\end{theorem}

Both bounds are achieved with equality on the shadowing orbit
$n = 2^L-1$, $T = L$ (\code{terras\_growth\_bound\_sharp}), where both
sides reduce to Lemma~\ref{lem:shadow}; neither can be improved.
The upper bound has the following striking consequence. Call density
$j/T$ \emph{supercritical} when $3^j > 2^{T-1}$, i.e.\ essentially
$j/T > \log 2/\log 3 \approx 0.6309$.

\begin{theorem}[density dichotomy; \code{no\_descent\_dichotomy}]
\label{thm:dichotomy}
If $n \le T^T(n)$ then, with $j = \jodd{T}{n}$, either
\[
2^T < 2 \cdot 3^{j}
\qquad\text{or}\qquad
n < 2^{\,T - j}.
\]
\end{theorem}

In words: an orbit segment that avoids descent either runs
supercritical odd density, or has already spent more than $\log_2 n$
halvings. Since each halving consumes one bit of $n$ --- a heuristic
made precise by the parity bijection of Section~\ref{sec:parity} ---
the second branch says the segment has exhausted the information
content of its starting point. In particular
(\code{early\_no\_descent\_forces\_density}), within the first
$\log_2 n$ worth of halvings, survival \emph{forces} supercritical
density.

Finally, the corridor matters because descent is the whole conjecture:

\begin{theorem}[descent reduction; \code{collatz\_iff\_descent}]
\label{thm:reduction}
The Collatz conjecture holds if and only if every $n \ge 2$ satisfies
$T^t(n) < n$ for some $t$.
\end{theorem}

The forward direction tracks the Terras orbit inside the Collatz orbit;
the converse is strong induction, descending through the first
$t$ with $T^t(n) < n$. Both directions are formal, so any future
descent theorem converts mechanically into the full conjecture.

\section{The critical line is sharp}\label{sec:critical}

Theorem~\ref{thm:dichotomy} forbids sub-critical survival inside the
information budget. The lower master inequality provides the exact
converse: supercritical \emph{prefixes} force survival.

\begin{theorem}[\code{no\_descent\_of\_supercritical\_prefix}]
\label{thm:prefix}
If $2^t \le 3^{\jodd{t}{n}}$ for all $t \le T$, then
$n \le T^t(n)$ for all $t \le T$.
\end{theorem}

\begin{proof}
By Theorem~\ref{thm:master} (lower bound) and $2^{j} \le 2^t \le 3^j$:
$2^t\,T^t(n) \ge 3^j(n+1) - 2^j \ge 2^t n + (3^j - 2^j) \ge 2^t n$.
\end{proof}

Which prefixes occur? By the realization theorem of
Section~\ref{sec:parity} (\code{parity\_pattern\_realized}), \emph{all}
of them: in particular the extremal word $1^j 0^s$ ($j$ odd steps, then
$s$ even ones), whose every prefix is supercritical as long as
$2^{j+s} \le 3^j$. This yields adversaries at every supercritical
density --- the Mersenne family of Corollary~\ref{cor:unbounded} is the
case $s = 0$, density $1$ --- and, choosing $j$ minimal, adversaries
hugging the critical line:

\begin{theorem}[\code{critical\_adversary}, \code{critical\_line\_sharp}]
\label{thm:sharp}
\emph{(i)} For all $j, s, N_0$ with $2^{j+s} \le 3^j$ there exists
$n \ge N_0$ with $\jodd{j+s}{n} = j$ and $n \le T^t(n)$ for all
$t \le j+s$.
\emph{(ii)} For every window $B$ and every $N_0$ there exists
$n \ge N_0$ with $n \le T^t(n)$ for all $t \le B$ and
\[
2^B \;\le\; 3^{\,\jodd{B}{n}} \;\le\; 3 \cdot 2^B .
\]
\end{theorem}

Part (ii) pins the odd density of an actual non-descending integer
within $O(1/B)$ of $\log 2/\log 3$ from above, for every $B$; part (i)
of Theorem~\ref{thm:dichotomy} forbids anything below. The
\emph{non-descent profiles achievable over a window $B$ are exactly a
strip of width $O(1/B)$ above the critical line, inhabited at every
scale by arbitrarily large integers}. By Theorem~\ref{thm:reduction},
the conjecture is the assertion that no single integer inhabits the
strip at every scale simultaneously. (Numerically, the construction at
$j = 100$ produces integers of density $0.6329$, against the critical
$0.6309$, with no descent through all $158$ steps.)

\section{The parity bijection and the binomial law}\label{sec:parity}

The constructions above rest on Terras' structural theorem
\cite{terras}, which we formalize in both directions. Congruence is
$\bmod\ 2^k$ throughout.

\begin{theorem}[\code{parity\_of\_modEq}, \code{modEq\_of\_parity}]
\label{thm:bijection}
For all $k$: $n \equiv m \pmod{2^k}$ if and only if
$T^t(n) \equiv T^t(m) \pmod 2$ for all $t < k$. Consequently the map
sending $r < 2^k$ to its parity vector
$\bigl(T^t(r) \bmod 2\bigr)_{t<k}$ is a bijection onto $\{0,1\}^k$.
\end{theorem}

The proof in both directions rides on the doubling identities of
Section~\ref{sec:corridor}: doubling carries congruence modulo $2^k$
to congruence modulo $2^{k+1}$ and back, and $3$ is invertible modulo
every power of $2$. Three corollaries do real work:

\begin{corollary}\label{cor:parity-consequences}
\emph{(i)} (\code{oddSteps\_modEq}) $\jodd{k}{\cdot}$ is a function of
$n \bmod 2^k$.
\emph{(ii)} (\code{parity\_pattern\_realized}) Every subset of
$\{0, \dots, k-1\}$ is the odd-step set of exactly one $r < 2^k$ ---
every finite behavior occurs.
\emph{(iii)} (\code{card\_oddSteps}, the binomial law)
\[
\#\bigl\{\, r < 2^k : \jodd{k}{r} = j \,\bigr\} = \binom{k}{j}.
\]
\end{corollary}

Item (ii) explains the failure of worst-case certificates
structurally: the adversarial pattern (all ones) is always present,
exactly once, at every scale --- and item (iii) says it and its
supercritical neighbours are binomially rare, which is the engine of
the next section.

\section{Almost every integer descends: Terras' theorem, finite form}
\label{sec:terras}

Fix $k$ and consider the dyadic window $[2^k, 2^{k+1})$, which
contains exactly one integer from each residue class modulo $2^k$. If
$n$ in the window does not descend within $k$ steps, the dichotomy
(with the window position supplying $n \ge 2^k \ge 2^{k-j}$) forces
supercritical density $2^k < 2\cdot 3^{\jodd{k}{n}}$; by
Corollary~\ref{cor:parity-consequences}(i) and (iii) the eligible
residues number $\sum_{j \,:\, 2^k < 2\cdot 3^j} \binom{k}{j}$. Two
elementary integer facts finish:
\[
2^m \sum_{\substack{j \le k \\ j \ge m}} \binom{k}{j}
\;\le\; \sum_{j \le k} \binom{k}{j}\, 2^j \;=\; (1+2)^k = 3^k
\qquad\text{(\code{choose\_tail\_bound})},
\]
and $2^k < 2 \cdot 3^j$ implies $j \ge \lfloor 17(k-1)/27 \rfloor$,
because $3^{17} \le 2^{27}$ (\code{tail\_exponent}; note
$17/27 = 0.6296\ldots$ approximates $\log 2/\log 3 = 0.6309\ldots$
from below).

\begin{theorem}[Terras' theorem, finite form;
\code{no\_descent\_window\_bound}]\label{thm:terras}
For every $k$,
\[
2^{\lfloor 17(k-1)/27 \rfloor} \cdot
\#\bigl\{\, n \in [2^k, 2^{k+1}) : n \le T^t(n) \text{ for all } t \le k
\,\bigr\} \;\le\; 3^k .
\]
\end{theorem}

The count is therefore at most $3^k/2^{\lfloor 17(k-1)/27\rfloor}
\approx 2^{0.955k}$: a fraction $\approx 2^{-0.045 k}$ of the window,
vanishing exponentially. The same bound holds for integers whose orbit
\emph{never} descends (\code{divergent\_window\_bound}): almost every
integer has finite stopping time. The formal proof uses no real
numbers, no measure theory, and no entropy estimates --- only the
weighted binomial theorem and one numeric inequality --- which we
found necessary as well as pleasant: the natural analytic route
through the popcount of $3^j$ runs into a genuinely open problem
(Section~\ref{sec:discussion}). Numerically the truth tracks the bound
comfortably: the actual fraction of non-descenders falls $0.50, 0.19,
0.063, 0.029$ at $k = 1, 4, 10, 18$.

Together with Section~\ref{sec:critical}: the strip above the critical
line is inhabited at every scale (Theorem~\ref{thm:sharp}), but its
population is an exponentially vanishing minority
(Theorem~\ref{thm:terras}). ``Almost all'' is now formal; ``all'' is
the conjecture.

\section{Cycle bounds}\label{sec:cycles}

On a cycle the two master inequalities pinch against each other.
Suppose $T^T(n) = n$ with $T \ge 1$, $n \ge 1$, and write
$j = \jodd{T}{n}$.

\begin{theorem}[\code{cycle\_has\_odd\_step},
\code{cycle\_three\_pow\_lt}, \code{cycle\_min\_bound}]
\label{thm:cycles}
$j \ge 1$; moreover $3^j < 2^T$ strictly, and
\[
n \cdot 2^T \;<\; n \cdot 3^j + 3^j\, 2^{\,T-j},
\qquad\text{i.e.}\qquad
n \;<\; \frac{3^j\, 2^{T-j}}{2^T - 3^j}.
\]
\end{theorem}

Every cycle element is thus bounded by how well $3^j$ approximates
$2^T$ from below --- the classical constraint
\cite{eliahou, simons-deweger}, now fully formal. The bound has teeth
at small $T$ because good rational approximations to $\log_2 3$ are
sparse. A finite computation --- performed by the Lean kernel itself
via plain \code{decide}, with no \code{native\_decide} and no trusted
code; the resulting lemma \code{cycle\_check\_183} depends on the
axiom \code{propext} alone --- verifies that for every $T \le 183$ and
every admissible $j$,
\[
3^j\, 2^{T-j} \;\le\; (2^T - 3^j)\cdot 2^{71},
\]
so every element of such a cycle is below $2^{71}$.

\begin{theorem}[\code{no\_small\_terras\_cycle}]\label{thm:nocycles}
Assume every $0 < m < 2^{71}$ reaches $1$ under the Collatz map (the
computational verification \cite{barina}, stated as a hypothesis).
Then every $n \ge 1$ with $T^T(n) = n$ for some $1 \le T \le 183$
reaches $1$; equivalently, no nontrivial cycle has $183$ or fewer
halvings.
\end{theorem}

The constant is sharp for this method at this verification bound: at
$(T, j) = (184, 116)$ the element bound first clears the verified
range ($3^{116} 2^{68}/(2^{184} - 3^{116}) \approx 2^{71.25}$), the
shadow of the approximation $116 \cdot \log_2 3 \approx 183.86$.
Pushing further requires a larger verified range or convergent-by-%
convergent Diophantine input \cite{eliahou}, not new ideas --- which
is itself consistent with Theorem~\ref{thm:nogo}: cycle exclusion is
finite-information reasoning, and provably cannot close the
conjecture.

\section{Certificate searches: from residue classes to digit automata}
\label{sec:search}

\subsection{The residue-class search}

Theorem~\ref{thm:nogo} was found by attempting its negation. For state
spaces $Q = \Z/2^a \times \Z/3^b$ one can search for additive
certificates $\varphi \colon Q \to \mathbb{R}$ mechanically. Model the
dynamics as a finite graph on $Q$: from a state $(r_2, r_3)$, the next
$T$-iterate's class is determined except for one bit --- each halving
consumes one bit of $2$-adic information --- so each state carries two
out-edges (the $3$-adic part is exact, $2$ being invertible modulo
$3^b$), with weight $-\log 2$ on even states and $+\log\frac32$ on odd
states. A certificate exists if and only if the associated linear
program is feasible, equivalently (by LP duality) if and only if every
cycle of this graph has negative mean weight, which Karp's algorithm
\cite{karp} decides exactly.

The graph over-approximates the true dynamics (both lifts of the lost
bit are kept), which is sound for worst-case certificates; and the
blocking cycle found below is realized by genuine integers, so no slack
is lost. The results:

\begin{center}
\begin{tabular}{cccl}
\toprule
$(a, b)$ & states & max mean cycle weight & verdict\\
\midrule
$(4, 0)$ & $16$  & $+0.405465$ & infeasible\\
$(6, 1)$ & $192$ & $+0.405465$ & infeasible\\
$(6, 2)$ & $576$ & $+0.405465$ & infeasible\\
$(5, 3)$ & $864$ & $+0.405465$ & infeasible\\
$(8, 1)$ & $768$ & $+0.405465$ & infeasible\\
\bottomrule
\end{tabular}
\end{center}

\noindent
The maximum mean cycle weight equals $\log\frac32 = 0.405465\ldots$
exactly, at every modulus, and the optimal cycle is always the same: the
self-loop at $n \equiv -1 \pmod{2^a}$. This is the $2$-adic fixed point
of Section~\ref{sec:shadow} seen at finite resolution, and tracing it
back to the integers $2^L - 1$ converts the failed search into
Theorem~\ref{thm:nogo}. The search tool
(\code{analysis/certificate\_search.py}) remains useful as a mechanical
refuter: any future claimed finite-modulus certificate must, by the
theorem, contain an error that the tool will locate as a positive-mean
cycle.

\subsection{The adversary spectrum: positive-drift rational cycles}
\label{sec:spectrum}

The Terras map preserves odd denominators: $T(a/q)$ is $(a/2)/q$ or
$((3a+q)/2)/q$. Every rational $a/q \in \Ztwo$ is shadowed by integers
$n \equiv a q^{-1} \pmod{2^L}$ for $\sim L$ steps, so every $T$-cycle
of rationals with \emph{positive drift} ($3^j > 2^P$ over the cycle)
is an adversary family in the sense of Section~\ref{sec:nogo}.
Enumerating all cycles with $q \le 99$ odd and starting numerators
$|a| \le 6q$ (\code{analysis/rational\_cycles.py}) finds $16$
positive-drift cycles in that range, all on negative rationals; the
strongest:

\begin{center}
\begin{tabular}{lcccc}
\toprule
cycle through & period $P$ & odd steps $j$ & density & drift ($\log_2$/step)\\
\midrule
$-1$      & $1$  & $1$ & $1.000$ & $+0.585$\\
$-65/49$  & $5$  & $4$ & $0.800$ & $+0.268$\\
$-19/11$  & $4$  & $3$ & $0.750$ & $+0.189$\\
$-5$      & $3$  & $2$ & $0.667$ & $+0.057$\\
$-17$     & $11$ & $7$ & $0.636$ & $+0.009$\\
\bottomrule
\end{tabular}
\end{center}

\noindent
The no-go theorem's $-1$ is the extreme point of a structured
spectrum whose densities descend toward the critical line --- the
dynamical counterpart of Theorem~\ref{thm:sharp}.

\subsection{Beyond bounded memory: digit automata and three killers}
\label{sec:automata}

By Remark~\ref{rem:unbounded}, certificates with \emph{unbounded} $W$
are not excluded, provided they carry a validity floor
$V(n) \ge n^{\delta}$. The natural such family: weighted automata
reading all binary digits of $n$ (least-significant first), $\log_2
V(n) = \log_2 n + \sum_i g(q_i, b_i)$, which can vary over
$n^{\Theta(1)}$ scales. We searched this family systematically
(\code{analysis/automaton\_potentials.py},
\code{analysis/transducer\_search.py}); since $\log V$ is linear in
the weight vector $g$, fitting $g$ against an adversary pool is a
piecewise-linear min--max problem, and we run full
counterexample-guided synthesis: fit, attack (Mersenne family,
rational-cycle shadows, exact minimal-weight strings by dynamic
programming, shift injections, hill-climbing over bit flips), add the
violator to the pool, repeat.

\emph{One-state results.} The popcount potential
$V(n) = n \cdot c^{\mathrm{ones}(n)}$ contracts along the Mersenne
orbit precisely when $c > 9/4$ --- the measured boundary
$\log_2 c = 1.17$ matches the prediction from
$\mathrm{ones}(3^j 2^{L-j}-1) \approx L - j/2$ --- so it kills the
no-go adversary. It dies elsewhere: the full two-parameter family
$V(n) = n\, 2^{\,g_1 \mathrm{ones}(n) + g_0 \mathrm{zeros}(n)}$ is
refuted across the entire $(g_0, g_1)$ plane.

\emph{Two-state results.} All $16$ two-state transition structures,
with the validity floor enforced as linear constraints (every cycle of
the weight graph must have mean $\ge -(1-\delta)$; without it the
optimizer returns inverted potentials $V \sim n^{-\gamma}$ that
certify nothing), are refuted by the CEGIS loop. The most instructive
candidate is the \emph{trailing-run automaton}, whose state records
whether the read head is still inside the initial run of ones --- so
it computes $\nu_2(n+1)$, the exact predictor of Mersenne rising
stretches --- with fitted potential $V(n) \approx n^{1.79} \cdot
4.6^{\,\nu_2(n+1)}$. It \emph{pre-pays} for every rise and genuinely
contracts along the entire Mersenne family. It dies in one step: for
$n = 2(2^{40}-1)$, the first halving re-aligns the read head with the
giant run of ones and releases the stored weight at once (measured
$+86.8$ in $\log_2 V$; the $30$-step window minimum stays at
$+52.9$).

Three refutation mechanisms, each with constructive adversaries,
emerged:
\begin{enumerate}
\item \textbf{Injection}: strings of extreme low weight (computable
exactly by dynamic programming on the automaton) whose statistics the
dynamics refills --- the orbit may even \emph{descend} while the
potential explodes (a measured case: $\Delta\log_2 n = -8.3$ over $40$
steps while $\Delta \log_2 V = +56$, the ones-density recovering from
$0.025$ toward typical).
\item \textbf{Global rewrite}: the odd step multiplies the entire
remaining bit string by $3$; digit statistics are not locally
controlled, and even positive-drift cycle shadows pump popcount-style
weights.
\item \textbf{Shift re-alignment}: halvings shift the tape, so any
LSB-anchored state has its stored potential released by a single even
step, as above.
\end{enumerate}
Mechanism 3 is evaded exactly by \emph{shift-coherent} weights ---
sliding $k$-gram statistics --- whose $k = 1, 2$ cases are killed by
mechanisms 1--2. We record the resulting conjecture:

\begin{conjecture}\label{conj:kgram}
No $k$-gram digit potential (for any $k$) with validity floor
$V \ge n^{\delta}$, window $B$, and factor $\rho < 1$ is a descent
certificate for $T$.
\end{conjecture}

Base-$2$ statistics die on $\times 3$; base-$3$ statistics die,
symmetrically, on halving. A weight statistic compatible with both
must live on the joint $2$-adic/$3$-adic structure --- the
$\times 2,\times 3$ rigidity landscape of Furstenberg's times-two,
times-three problem --- which we take as quantitative evidence that
the difficulty of the conjecture is located exactly where the folklore
places it. These are experimental refutations: every kill is an
explicit, checkable violator, but ``no certificate exists in this
family'' remains conjecture pending a proof of mechanism 1 or 3 in the
style of Section~\ref{sec:critical} (prescribed parity words rather
than explicit bit strings; see Section~\ref{sec:discussion}).

\section{Case study: how claims like Corollary~\ref{cor:axiom}'s
negation get ``verified''}
\label{sec:casestudy}

Arguments of the genre formalized in Definition~\ref{def:cert}
typically proceed by tracking, for each odd residue class modulo
$2^{14}$ (say), an affine form $c \cdot n + d$ along the orbit, with
worst-case handling of boundary cases, and reporting that every class
descends within some bound --- $53$ steps, for the claim refuted in
Corollary~\ref{cor:axiom}. We implemented exactly such a tracker to
locate the characteristic flaw. It is this: the class of $n$ modulo
$2^{14}$ determines the orbit's parity decisions only until the
cumulative number of halvings reaches $14$ --- each halving spends one
bit, exactly as in Theorem~\ref{thm:bijection} --- while the tracker
runs for up to $53$ steps. Beyond the budget, the tracked trajectory
silently follows the class representative rather than the class. For
the class of $n = 10087$, the tracker first diverges from reality at
step $12$ (cumulative halvings $15$), then certifies a descent that
$n = 10087$ itself violates: its actual stopping time is $66$ Terras
steps.

We record this because it is the \emph{typical} error of the genre,
and because Theorem~\ref{thm:nogo} explains its inevitability: had the
precision accounting been done correctly at modulus $2^{14}$, the
verification would have failed on the classes near $-1$; enlarging the
modulus relocates but never removes them. By the binomial law
(Corollary~\ref{cor:parity-consequences}) the classes with no
guaranteed early descent number $\sum_{j > 0.63k}\binom{k}{j} \sim
2^{0.95k}$ at modulus $2^k$: their density tends to $0$ --- that is
Theorem~\ref{thm:terras} --- but their count grows exponentially,
which is why ``almost every'' is where this method provably stops.

\section{The formalization}\label{sec:formal}

All theorems of Sections \ref{sec:shadow}--\ref{sec:cycles} are
formalized in Lean~4 over \textsf{mathlib} \cite{mathlib}, about
$2{,}100$ lines across seven files:

\begin{center}
\begin{tabular}{lll}
\toprule
file & contents & headline results\\
\midrule
\code{Basic.lean}    & maps, $2$-adic valuation & \code{CollatzConjecture}\\
\code{NoGo.lean}     & \S\ref{sec:shadow}--\ref{sec:nogo} & \code{no\_finite\_certificate}\\
\code{Density.lean}  & \S\ref{sec:corridor} & \code{no\_descent\_dichotomy}, \code{collatz\_iff\_descent}\\
\code{Cycles.lean}   & \S\ref{sec:cycles} & \code{cycle\_min\_bound}, \code{no\_small\_terras\_cycle}\\
\code{Parity.lean}   & \S\ref{sec:parity} & \code{card\_oddSteps}, \code{parity\_pattern\_realized}\\
\code{Terras.lean}   & \S\ref{sec:terras} & \code{no\_descent\_window\_bound}\\
\code{Critical.lean} & \S\ref{sec:critical} & \code{critical\_line\_sharp}\\
\bottomrule
\end{tabular}
\end{center}

The no-go statement, as type-checked:

\begin{leanblock}
theorem no\_finite\_certificate\\
\hspace*{1.5em}(W : $\N$ $\to$ $\Q$) (lo hi : $\Q$) (hlo : 0 < lo)\\
\hspace*{1.5em}(hW : $\forall$ n, lo $\le$ W n $\wedge$ W n $\le$ hi)\\
\hspace*{1.5em}(B N$_0$ : $\N$) ($\rho$ : $\Q$) (h$\rho_0$ : 0 $\le$ $\rho$) (h$\rho$ : $\rho$ < 1)\\
\hspace*{1.5em}(hcert : $\forall$ n, N$_0$ $\le$ n $\to$ $\exists$ t, 1 $\le$ t $\wedge$ t $\le$ B $\wedge$\\
\hspace*{3em}(terras\_iter t n : $\Q$) * W (terras\_iter t n)\\
\hspace*{4.5em}$\le$ (n : $\Q$) * W n * $\rho$) :\\
\hspace*{1.5em}False
\end{leanblock}

\noindent and Terras' theorem in finite form:

\begin{leanblock}
theorem no\_descent\_window\_bound (k : $\N$) :\\
\hspace*{1.5em}2 \^{} (17 * (k - 1) / 27) *\\
\hspace*{3em}((Finset.Ico (2 \^{} k) (2 \^{} (k + 1))).filter\\
\hspace*{4.5em}(fun n => $\forall$ t, t $\le$ k $\to$ n $\le$ terras\_iter t n)).card\\
\hspace*{3em}$\le$ 3 \^{} k
\end{leanblock}

\paragraph{Audit.} Beyond a successful clean build, we verified:
(i) \code{\#print axioms} on the $24$ headline theorems reports only
\code{propext}, \code{Classical.choice}, and \code{Quot.sound} ---
Lean's standard foundations; the kernel computation
\code{cycle\_check\_183} depends on \code{propext} alone;
(ii) the sources contain no \code{sorry}, no \code{axiom}
declarations, no \code{unsafe} or \code{partial} definitions, no
\code{@[implemented\_by]} and no \code{native\_decide}, so no result
depends on compiled code or escape hatches; (iii) every quantitative
statement was cross-checked against independent Python
implementations: the master inequalities on $20{,}000$ random
$(n, T)$, the dichotomy on all non-descending cases among them, the
binomial law exhaustively for $k \le 12$, the window bound and actual
survivor counts for $k \le 18$, the critical adversaries to $j = 100$,
the sharpness identity at $L \in \{5, 10, 20, 28\}$, and the
$528{,}958{,}107{,}647 = T^{17}(2^{29}-1)$ spot check of
Lemma~\ref{lem:shadow}; and (iv) the only hypothesis anywhere that is
not formally proved is the explicitly displayed verification
hypothesis of Theorem~\ref{thm:nocycles}.

We are aware of formalizations of Collatz \emph{computations} and of
the conjecture's statement in several proof assistants; we are not
aware of a prior machine-verified impossibility theorem for a class of
Collatz proof strategies, nor of a prior machine-verified density
result for the map (Theorem~\ref{thm:terras}), though we would welcome
correction.

\section{What remains}\label{sec:discussion}

The formal results fence off worst-case bounded-memory arguments
(Theorem~\ref{thm:nogo}), and the experiments of
Section~\ref{sec:automata} push the fence, so far without exception,
through the unbounded digit-automaton families. What remains
available:

\begin{enumerate}
\item \textbf{Density and averaging arguments.} These carry unbounded
information and are untouched; indeed Section~\ref{sec:terras} is one,
now formal, and Tao's result that almost all orbits attain almost
bounded values \cite{tao} is their modern summit. But they prove
``almost all,'' and Theorem~\ref{thm:sharp} shows the exceptional
strip they permit is genuinely inhabited at every finite scale.
\item \textbf{Cycle exclusions} (Section~\ref{sec:cycles}): real
theorems, conditional only on verified ranges, but provably
finite-information --- they bound cycle lengths, never settle
divergence.
\item \textbf{Joint $2$-adic/$3$-adic structure.} The three killers of
Section~\ref{sec:automata} leave exactly one road into certificate
land: statistics quasi-invariant under both $\times 3$-with-carry and
halving, i.e.\ living on $\Ztwo \times \Z_3$. This is the
times-two/times-three rigidity landscape, and we read the experiments
as quantitative evidence that the conjecture's difficulty is exactly
the difficulty of that landscape.
\item \textbf{Genuinely new ideas} exploiting the specific arithmetic
of $2$ and $3$. Some constraint of this kind is mandatory: the
generalized Collatz problem is undecidable \cite{conway, kurtz-simon},
so no argument indifferent to the particular map can succeed.
\end{enumerate}

Two concrete formal targets suggest themselves. First,
Conjecture~\ref{conj:kgram}: the injection mechanism is constructive,
and the realization machinery of Section~\ref{sec:parity} can
prescribe parity words without controlling explicit bit strings --- a
route that matters because the naive analysis of even the popcount
potential on the Mersenne orbit needs the popcount of $3^j$, on which
nothing beyond Stewart-type logarithmic bounds is known
\cite{stewart}; any proof must go around that open problem, and our
finite-form Terras proof shows such detours exist. Second, the
quantitative no-go: certificates with window $B(n) = o(\log n)$ should
be excludable by the chaining of Theorem~\ref{thm:nogo} run inside
Theorem~\ref{thm:sharp}'s adversaries.

The conjecture (descent part) is equivalent, with sharp constants on
both sides now formal, to the assertion that no orbit tracks the
critical strip $2^T \le 3^{\,\jodd{T}{n}} \le 3 \cdot 2^T$ at every
scale $T$. Random parities give density $\frac12$, far below critical;
the binomial law makes the strip exponentially thin; the realization
theorem keeps it nonempty at every scale; and undecidability of the
generalized problem warns that only the specific arithmetic of $2$ and
$3$ can finish. We see this paper as a fence around the parts of the
search space that are provably or demonstrably empty --- with the
gate, the critical line, now surveyed exactly --- and as a case study
in making impossibility \emph{informative}.

\subsection*{Acknowledgments}
The formalization, the searches, and this text were produced in
collaboration with Claude (Anthropic). The Lean community's
\textsf{mathlib} made the formalization routine.

\subsection*{Reproducibility}
All Lean sources, the analysis scripts (deterministic seeds), and
build instructions are in the repository accompanying this paper.
\code{lake build} checks every theorem; \code{\#print axioms}
reproduces the audit; each script regenerates the corresponding table
or experiment of Sections \ref{sec:search} in minutes on commodity
hardware.

\begin{thebibliography}{10}

\bibitem{barina} D.~Barina, Convergence verification of the Collatz
problem, \emph{The Journal of Supercomputing} \textbf{77} (2021),
2681--2688.

\bibitem{conway} J.~H.~Conway, Unpredictable iterations, in
\emph{Proceedings of the 1972 Number Theory Conference}, University of
Colorado, Boulder, 1972, pp.~49--52.

\bibitem{eliahou} S.~Eliahou, The $3x+1$ problem: new lower bounds on
nontrivial cycle lengths, \emph{Discrete Mathematics} \textbf{118}
(1993), 45--56.

\bibitem{karp} R.~M.~Karp, A characterization of the minimum cycle mean
in a digraph, \emph{Discrete Mathematics} \textbf{23} (1978), 309--311.

\bibitem{krasikov-lagarias} I.~Krasikov and J.~C.~Lagarias, Bounds for
the 3x+1 problem using difference inequalities, \emph{Acta Arithmetica}
\textbf{109} (2003), 237--258.

\bibitem{kurtz-simon} S.~A.~Kurtz and J.~Simon, The undecidability of
the generalized Collatz problem, in \emph{Theory and Applications of
Models of Computation (TAMC 2007)}, Lecture Notes in Computer Science
4484, Springer, 2007, pp.~542--553.

\bibitem{lagarias-survey} J.~C.~Lagarias, The $3x+1$ problem and its
generalizations, \emph{American Mathematical Monthly} \textbf{92}
(1985), 3--23.

\bibitem{lagarias-book} J.~C.~Lagarias (ed.), \emph{The Ultimate
Challenge: The $3x+1$ Problem}, American Mathematical Society, 2010.

\bibitem{mathlib} The mathlib Community, The Lean mathematical library,
in \emph{Proceedings of the 9th ACM SIGPLAN International Conference on
Certified Programs and Proofs (CPP 2020)}, ACM, 2020, pp.~367--381.

\bibitem{simons-deweger} J.~Simons and B.~de~Weger, Theoretical and
computational bounds for $m$-cycles of the $3n+1$ problem, \emph{Acta
Arithmetica} \textbf{117} (2005), 51--70.

\bibitem{stewart} C.~L.~Stewart, On the representation of an integer
in two different bases, \emph{Journal f\"ur die reine und angewandte
Mathematik} \textbf{319} (1980), 63--72.

\bibitem{tao} T.~Tao, Almost all orbits of the Collatz map attain
almost bounded values, \emph{Forum of Mathematics, Pi} \textbf{10}
(2022), e12.

\bibitem{terras} R.~Terras, A stopping time problem on the positive
integers, \emph{Acta Arithmetica} \textbf{30} (1976), 241--252.

\end{thebibliography}

\end{document}
